\subsection{The sphere criterion and the general-localizer answer to C048}

\begin{proposition}[Known: the localizer generated by a sphere]\label{F:sphere-localizer}
Let $W$ be any fundamental localizer, with no accessibility assumption,
and let $n\geq0$. The following conditions are equivalent:
\begin{enumerate}
\item the category $\Delta/\partial\Delta^{n+2}$ is $W$-aspherical;
\item $W_n\subseteq W$.
\end{enumerate}
In particular, $S^2$ is $W$-aspherical if and only if $W_1\subseteq W$.
Here an assertion about the asphericity of a space means the
corresponding assertion for a category having its weak homotopy type.
\end{proposition}
\begin{proof}
This is \citet[Proposition~9.2.17, equivalence (a)--(b), pp.~369--370]{Cisinski2006Prefaisceaux}.
The same source's Corollary~9.2.15 identifies its $n$-equivalences with
isomorphisms on $\pi_0$ and on all based $\pi_i$ for $1\leq i\leq n$,
so its convention agrees with Definition~\ref{D:n-convention}.
The proposition is stated for an arbitrary fundamental localizer.
Its hypotheses do not require that $W$ itself be accessible.
The simplex category $\Delta/\partial\Delta^{n+2}$ has the weak
homotopy type of the sphere $S^{n+1}$. Since
$W_\infty\subseteq W$, replacing this category by another representative
of the same weak homotopy type does not affect its $W$-asphericity.
\end{proof}

\begin{lemma}[Proved: the groupoid comparison for an ordinary pushout]\label{F:span-groupoid}
Let $i:A\to B$ be injective on objects, and let $j:A\to C$ be any functor.
Let $\mathcal G$ be the Grothendieck construction of the covariant
diagram $B\leftarrow A\to C$ indexed by the three-object outgoing
span. The canonical comparison
\[
 q:\mathcal G\longrightarrow P:=B\amalg_A C
\]
belongs to $W_1$.
\end{lemma}
\begin{proof}
Write $\Pi E$ for the groupoid completion of a category $E$.
Its components and automorphism groups are respectively the components
and fundamental groups of $BE$, by the edge-path presentation of the
nerve. It therefore suffices to prove that
$\Pi q:Q:=\Pi\mathcal G\to\Pi P$ is an equivalence of groupoids.

The category $\mathcal G$ has tagged copies $a_A,b_B,c_C$ of the objects
of $A,B,C$, and canonical arrows
\[
 u_a:a_A\longrightarrow i(a)_B,\qquad
 v_a:a_A\longrightarrow j(a)_C.
\]
Together with the arrows within the three copies, these generate all
its arrows, with the naturality relations
\[
 u_{a'}\alpha_A=i(\alpha)_B u_a,\qquad
 v_{a'}\alpha_A=j(\alpha)_C v_a
 \quad(\alpha:a\to a').
\]
In $Q$ all these arrows are invertible.

Define a functor $B':B\to Q$ on objects by
\[
 B'(b)=
 \begin{cases}
 j(a)_C,&b=i(a),\\
 b_B,&b\notin i(\operatorname{Ob}A).
 \end{cases}
\]
This is unambiguous because $i$ is injective on objects. Put
$\theta_{i(a)}=v_a u_a^{-1}:i(a)_B\to j(a)_C$ and
$\theta_b=\id_{b_B}$ off the image of $i$, and define
$B'(\beta)=\theta_{b'}\beta_B\theta_b^{-1}$ for $\beta:b\to b'$.
The displayed naturality relations give the strict equality of
functors $B'i=C_{\!C}j$, where $C_{\!C}:C\to Q$ is the tagged-copy
functor. The universal property of the ordinary categorical pushout,
followed by groupoid completion, supplies
$r:\Pi P\to Q$.

The equality $(\Pi q)r=\id_{\Pi P}$ follows on the generating copies
of $B$ and $C$, since every $\theta_b$ maps to an identity.
There is a natural isomorphism
$\id_Q\Rightarrow r(\Pi q)$ whose components are $\theta_b$ on
$b_B$, identities on $c_C$, and $v_a$ on $a_A$.
Naturality on $u_a$ is the equality $\theta_{i(a)}u_a=v_a$;
naturality on $v_a$ is immediate; the remaining naturality equations
are the conjugation formula and the two displayed relations.
This proves the required equivalence and hence $q\in W_1$.
\end{proof}

\begin{theorem}[Proved: all object-injective maps above the one-type localizer]\label{F:sphere-cofibrations}
If a fundamental localizer $W$ makes $S^2$ aspherical, then every
functor injective on objects is a $W$-cofibration in the manuscript's
sense.
\end{theorem}
\begin{proof}
Proposition~\ref{F:sphere-localizer} gives $W_1\subseteq W$.
Fix an object-injective $i:A\to B$, a functor $j:A\to C$, and a
$W$-equivalence $f:C\to C'$ under $A$. The comparison maps from
the two span Grothendieck constructions to
$B\amalg_A C$ and $B\amalg_A C'$ belong to $W_1$ by
Lemma~\ref{F:span-groupoid}, hence belong to $W$.

The functor between the two Grothendieck constructions also belongs
to $W$. For completeness, this uses only the relative Theorem~A
axiom for a fundamental localizer. For a covariant diagram
$D:J\to\Cat$ and $s\in J$, the slice of its Grothendieck construction
over $s$ retracts to $D(s)$ by
\[
 (t,x,\alpha:t\to s)\longmapsto D(\alpha)(x).
\]
Inclusion of the fibre is a section, and the arrows
$(\alpha,\id)$ give a natural transformation from the identity to
the inclusion followed by the retraction. Thus this is a
$W_\infty$-equivalence, naturally in the diagram.
An objectwise $W$-equivalence of diagrams consequently induces a
$W$-equivalence on each such slice, and the relative axiom gives
a $W$-equivalence on Grothendieck constructions.
Apply this to the objectwise map $(\id_B,\id_A,f)$ of spans.

The square of canonical comparisons commutes. Two-out-of-three
therefore gives
$B\amalg_A C\to B\amalg_A C'\in W$, exactly the required
preservation of $W$-equivalences by the pushout functor.
\end{proof}

\begin{corollary}[Proved: the complete C048 dichotomy]\label{F:C048-dichotomy}
Fix an alternating interval $\Phi_n$ with $n\geq2$ and its endpoint
sieve $j:\{0,2n\}\hookrightarrow\Phi_n$. For every fundamental
localizer $W$, the maximal class of full sieve inclusions satisfying
C048 is as follows:
\begin{enumerate}
\item If $S^2$ is $W$-aspherical, all full sieve inclusions form such a
class.
\item If $S^2$ is not $W$-aspherical, the only possible members are
isomorphisms, and the class of isomorphisms satisfies the requirements.
\end{enumerate}
No accessibility assumption on $W$ is needed.
\end{corollary}
\begin{proof}
In the first case Theorem~\ref{F:sphere-cofibrations} makes every
full sieve inclusion a $W$-cofibration. Full sieve inclusions are
preserved by arbitrary categorical cobase change. They are also
preserved by the displayed pushout-product operation: if
$U\hookrightarrow X$ is a full sieve, its pushout-product with $j$
is the inclusion of the full sieve
\[
 (U\times\Phi_n)\cup(X\times\{0,2n\})
 \ \hookrightarrow\ X\times\Phi_n.
\]
Thus the class of all full sieve inclusions has both requested
stability properties and is maximal.
The second case is Theorem~\ref{F:interval-impossibility}.
\end{proof}
